\documentclass[12pt]{article}
\usepackage[utf8]{inputenc}
\usepackage[T1]{fontenc}
\usepackage{hyperref}
\title{Yau Language Resources for K-12 Curriculum Development}
\author{Compiled by Claude Sonnet 5.5 (claude-sonnet-5-5) for Andr\'as Kornai}
\date{}
\begin{document}
\maketitle

\section{Introduction}
Yau (ISO 639-3: yyu; Glottocode yaus1235) is an endangered Torricelli language of the Wapei group. It is spoken in Senim village, East Wapei Rural LLG, Sandaun Province, Papua New Guinea. Estimates of speaker numbers differ: about 200 (2014, per Wikipedia) and about 140 (Verbix). It should not be confused with Yawa, a separate language family on Yapen Island in Indonesian Papua. The search found very little on Yau, and no teaching materials.

\section{Primary Educational Resources}
No primers, textbooks or school programmes for Yau turned up in the search.

\section{Linguistic and Reference Materials}
\begin{itemize}
\item Overview article on the language: \url{https://en.wikipedia.org/wiki/Yau_language_(Torricelli)}
\item Overview of the Wapei languages: \url{https://en.wikipedia.org/wiki/Wapei_languages}
\item Overview of the Torricelli family: \url{https://static.hlt.bme.hu/semantics/external/pages/jokuc/en.wikipedia.org/wiki/Torricelli_languages.html}
\end{itemize}
No grammar or dictionary of Yau appeared in the search results.

\section{Religious and Community Resources}
The search found no scripture translations or community organisations for Yau.

\section{Digital and Audio Resources}
\begin{itemize}
\item Endangered Languages Project entry: \url{https://www.endangeredlanguages.com/lang/3158?hl=pt-br}
\item Verbix endangered-languages entry: \url{https://docs.verbix.com/EndangeredLanguages/Yau}
\end{itemize}
These pages are catalogue entries only. I did not verify that they offer any audio or text.

\section{Curriculum Development Context}
Yau is still spoken by all adults but not all young people, and it is not known to be taught in schools. Documentation appears minimal. A curriculum would first require community-led documentation (a word list, an orthography and recordings) before any materials could be written. Nothing more reliable was found.
\end{document}